\section{Introduction}
An image communication system with a hard per-image symbol budget must decide which information to retain before transmission. The decision depends on image-dependent source-code length, channel protection, and reconstruction from the retained representation. A generative receiver can supply missing information, but its output must remain constrained by what was delivered. Pixel fidelity, perceptual similarity, feature consistency, and processing cost can favor different operating points.

Digital image coding with error correction and learned joint source--channel coding (JSCC) offer different approaches~\cite{deepjscc,swinjscc}. Truncation of progressive descriptions to fit a channel budget is a classical problem~\cite{chande}. Ada-TokenCom combines prefix entropy coding, tail generation, and joint prefix/MCS selection under an average resource constraint~\cite{adatokencom}; TONIC studies task-oriented token communication~\cite{tonic}. ARPC constructs progressive descriptions with an autoregressive prior~\cite{arpc}. We use the next-scale representation of visual autoregressive modeling (VAR)~\cite{var}: complete coarse maps describe the full feature grid at low resolution, finer maps add residual detail, and earlier scales condition later generation. This structure supports reconstruction from a coarse description with progressively more residual information.

Fixed frames do not generally align with complete-scale endpoints. Their token-count steps are large, while terminated entropy-code lengths vary across images. The next complete scale can exceed the frame even when the preceding endpoint leaves room. Within-scale endpoints retain part of this boundary map, filling the gap between whole scales. The next-scale factorization also lets both ends construct matching coding probabilities from the complete coarse maps, so the extra endpoints can be decoded without transmitting probability tables.

NeST-Com adapts coarse-to-fine token prefixes to fixed-budget digital transmission through within-scale endpoints, actual-length fitting, and prefix-constrained reconstruction. Received boundary positions update both the residual latent and the context used to generate finer scales, using VAR's known-token overwrite operation. With shared model components, quarter-scale endpoints improve PSNR by 0.32--0.70~dB over whole-scale transmission and raise entropy-stream utilization from 77.70\% to 95.21\% at 10~dB. Controlled endpoint and receiver interventions examine these mechanisms. Six-system comparisons at $N=1024$ and a focused, independently calibrated $N=512$ comparison locate the resulting pixel--perceptual tradeoffs under two symbol budgets.
